% =============================================================================
%  6-trabajo-futuro.tex — Importancia de la investigación y proyección doctoral
% =============================================================================

\section{Importancia de la Investigación y Proyección Doctoral}

\begin{frame}{Importancia de la investigación}
  \begin{enumerate}
    \item \textbf{Política pública.}
          La CDMX y el Estado de México cuentan con un marco cuantitativo para
          evaluar y planear conjuntamente el Anillo Periférico como eje de
          transporte masivo.

    \item \textbf{Metodología urbana.}
          La Clasificación Garibelt y el Espectro MJG ofrecen un diagnóstico
          multidimensional no compensable. Un mejor tiempo de viaje y un
          déficit de cobertura quedan visibles por separado.

    \item \textbf{Soberanía de datos.}
          El código auditable y las imágenes Docker permiten procesar datos GTFS
          y reproducir los indicadores sin licencias comerciales de software
          para el cálculo.
  \end{enumerate}

  \medskip
  {\footnotesize\itshape Fuente: conclusiones de la tesis, cap.~7.\par}
\end{frame}

\begin{frame}{Agenda doctoral: calibración y demanda}
  \begin{block}{Análisis de sensibilidad de la escala Garibelt}
    Evaluar cómo cambian las bandas de diagnóstico al variar los umbrales
    ($0.25$, $0.50$ y $0.75$), las cotas de tiempo (85--180~min) y los
    coeficientes de fricción vial ($CF_h$).
  \end{block}

  \begin{block}{Integración de topología y demanda}
    Acoplar el grafo ponderado con las matrices de la Encuesta
    Origen--Destino del INEGI y con volúmenes observados de pasajeros para
    evaluar los escenarios a partir de los desplazamientos y la afluencia.
  \end{block}
\end{frame}

\begin{frame}{Agenda doctoral: vulnerabilidad y aplicación territorial}
  \begin{block}{Cálculo formal de la vulnerabilidad ($\Delta E$)}
    Simular la remoción de Macro-Hubs completos y medir la caída de eficiencia
    global para comparar la resistencia de la red base y de los escenarios
    anillares ante fallas.
  \end{block}

  \begin{block}{Replicabilidad internacional y Anillo Periférico Exterior}
    Probar el método en otras metrópolis latinoamericanas con datos GTFS,
    como Bogotá y Santiago. Modelar un anillo exterior para estimar qué
    cobertura y conectividad podría aportar al cuarto contorno metropolitano.
  \end{block}
\end{frame}

% ── Bibliografía ─────────────────────────────────────────────────────────────
\section*{Referencias}
\begin{frame}[allowframebreaks]{Referencias}
  \bibliographystyle{apalike}
  \bibliography{referencias}
\end{frame}

% ── Cierre institucional ─────────────────────────────────────────────────────
\begin{frame}[plain]
  \begin{center}
    {\color{ColorAcento}\rule{0.6\textwidth}{1.5pt}}\par\medskip
    {\Large\bfseries\color{ColorPrincipal} Gracias\par}
    \medskip
    \begin{minipage}{0.73\textwidth}
      \centering
      \itshape La tinta escribrá las palabras que nuestras voces no fueron capaces de hablar.
    \end{minipage}\par
    \medskip
    {\normalsize \Autor\par}
    {\small \MateriaCorta\par}
    {\small \Semestre\par}
    \medskip
    {\color{ColorAcento}\rule{0.6\textwidth}{1.5pt}}\par
  \end{center}
\end{frame}
